\chapter{归一化、Softmax 与数值稳定性}

Transformer 推理里有许多看似“小”的算子：LayerNorm、RMSNorm、Softmax、激活函数、残差加法。它们的 FLOPs 可能远少于矩阵乘，但端到端性能中不能忽略，因为它们常常读写完整向量或矩阵，计算强度低，且位于关键路径上。更重要的是，这些算子对数值稳定性敏感，不能只为速度牺牲正确性。

\section{Softmax 为什么要减最大值}

Softmax 定义为 $\mathrm{softmax}(x_i)=\exp(x_i)/\sum_j \exp(x_j)$。如果直接对很大的 $x_i$ 求指数，可能溢出；如果很多值很小，可能下溢到零。常见稳定写法是先减去最大值 $m=\max_i x_i$，再计算 $\exp(x_i-m)$。这样最大指数为 1，其余不超过 1，不改变最终比例。

\begin{lstlisting}[language=C++,caption={稳定 Softmax：先减最大值，再归一化}]
#include <algorithm>
#include <cmath>
#include <cstdint>
#include <span>

void softmax_stable(std::span<const float> input,
                    std::span<float> output) {
    float max_value = input[0];
    for (std::int64_t i = 1; i < static_cast<std::int64_t>(input.size()); ++i) {
        max_value = std::max(max_value, input[static_cast<std::size_t>(i)]);
    }

    float sum = 0.0F;
    for (std::int64_t i = 0; i < static_cast<std::int64_t>(input.size()); ++i) {
        const float value =
            std::exp(input[static_cast<std::size_t>(i)] - max_value);
        output[static_cast<std::size_t>(i)] = value;
        sum += value;
    }

    const float inv_sum = 1.0F / sum;
    for (std::int64_t i = 0; i < static_cast<std::int64_t>(output.size()); ++i) {
        output[static_cast<std::size_t>(i)] *= inv_sum;
    }
}
\end{lstlisting}

这个实现有三次线性扫描：求最大值，求指数和总和，归一化输出。它清楚但不是最少内存流量。优化版本会尝试融合扫描、使用向量化指数近似、按块处理长行，或者在 Attention 中避免显式写出完整 Softmax 矩阵。但无论怎么优化，“减最大值”这个稳定性原则不能丢。

\section{归约的性能结构}

Softmax、LayerNorm 和 RMSNorm 都包含归约。归约的特点是多个输入合成少量统计量，例如最大值、总和、均值、方差或平方和。归约不只是数学操作，也是一种数据依赖结构。单线程里，归约容易形成累加依赖链；多线程里，需要把局部统计量合并；SIMD 里，需要 lane 内向量运算和 lane 间水平归约。

LayerNorm 需要均值和方差。朴素实现可能先算均值，再算方差，再归一化，至少三次扫描。若向量长度很大，内存流量明显；若向量长度很小，函数调用和调度开销可能显著。高性能实现会根据 hidden size、数据类型和平台选择不同策略。

\begin{lstlisting}[language=C++,caption={RMSNorm 的教学实现：平方和归约后缩放}]
#include <cmath>
#include <cstdint>
#include <span>

void rms_norm(std::span<const float> x,
              std::span<const float> weight,
              std::span<float> y,
              float epsilon) {
    float square_sum = 0.0F;
    for (std::int64_t i = 0; i < static_cast<std::int64_t>(x.size()); ++i) {
        const float value = x[static_cast<std::size_t>(i)];
        square_sum += value * value;
    }

    const float mean_square =
        square_sum / static_cast<float>(x.size());
    const float scale = 1.0F / std::sqrt(mean_square + epsilon);
    for (std::int64_t i = 0; i < static_cast<std::int64_t>(x.size()); ++i) {
        y[static_cast<std::size_t>(i)] =
            x[static_cast<std::size_t>(i)] * scale *
            weight[static_cast<std::size_t>(i)];
    }
}
\end{lstlisting}

\section{近似函数和误差边界}

指数、倒数平方根、GELU、SiLU 这类函数比加法和乘法昂贵。优化实现经常使用近似算法或硬件近似指令，然后通过误差分析保证模型质量不受明显影响。这里的关键是“误差边界”而不是“看起来差不多”。近似函数要说明最大误差、平均误差、输入范围和对后续归一化的影响。

AI 推理通常允许一定数值误差，因为训练和推理本身已经涉及浮点舍入、量化和不同后端差异。但误差不是无限的。Softmax 对极端 logits 敏感，归一化对 epsilon 和累加精度敏感，量化对离群值敏感。性能优化必须与正确性测试绑定。

\begin{warningbox}
不要为了少一次扫描就丢掉稳定性。Softmax 不减最大值、方差计算不处理 epsilon、低精度累加没有误差检查，都可能在普通样例上看不出问题，却在真实模型或长序列上产生严重错误。
\end{warningbox}

\section{融合机会}

归一化和逐元素激活常常适合融合。例如 RMSNorm 可以把缩放和权重乘法放在同一次输出扫描里；MLP 中的 bias、GELU 和乘法门控可以融合，减少中间张量。Softmax 在 Attention 中更特殊，如果显式产生 $QK^T$ 的完整矩阵，再对它做 Softmax，然后再乘 $V$，中间流量会非常大。FlashAttention 这类思想的核心，就是分块计算并在线维护 Softmax 统计量，避免完整中间矩阵落到慢内存。

\begin{keyidea}
Transformer 中的小算子往往不是“小问题”。它们的单次 FLOPs 不高，但访问频繁、处在关键路径、涉及归约和数值稳定性。优化时要同时看内存流量和误差边界。
\end{keyidea}
